% SEGMENT OLP-0043-S01
% Part: sets-functions-relations
% Chapter: arithmetization
% Section: rationals

\documentclass[../../../include/open-logic-section]{subfiles}

\begin{document}

\olfileid{sfr}{arith}{rat}
\olsection{$\Int$ இலிருந்து $\Rat$ க்கு}

முறைசாராத கணக்கோட்பாட்டைப் பயன்படுத்தி இயல் எண்களிலிருந்து
முழுக்களை எவ்வாறு கட்டமைப்பது என இப்போதுதான் கண்டோம்.
அதே போன்ற முறையில் முழுக்களிலிருந்து விகிதமுறு எண்களை எவ்வாறு
கட்டமைப்பது என இப்போது காண்போம். நமது தொடக்க உணர்தல்கள்:
\begin{enumerate}
	\item ஒவ்வொரு விகிதமுறு எண்ணையும் $\nicefrac{i}{j}$ என்ற
    வடிவில் எழுதலாம்; இங்கு $i$, $j$ இரண்டும் முழுக்கள்,
    ஆனால் $j$ பூச்சியமற்றது.
	\item $\nicefrac{i}{j}$ என்ற கோவையில் குறியாக்கப்பட்டுள்ள தகவலை,
    $\tuple{ i, j}$ என்ற வரிசைப்படுத்தப்பட்ட ஜோடியாலும்
    அதேபோலக் குறியாக்கலாம்.
\end{enumerate}
$\Int \times (\Int \setminus \{0_\Int\})$ இலிருந்து எடுக்கப்பட்ட
வரிசைப்படுத்தப்பட்ட ஜோடிகளாகவே விகிதமுறு எண்களைக் \emph{கருதுவது}
தெளிவான அணுகுமுறையாக இருக்கும்.
ஆனால் முன்புபோல, அது அளவுக்கு அதிகமாக முறைசாராதது;
ஏனெனில் $\nicefrac{3}{2} = \nicefrac{6}{4}$ என நமக்குத் தேவை,
ஆனால் $\tuple{ 3, 2}\neq \tuple{ 6, 4}$.
மேலும் பொதுவாக, பின்வருவது நமக்குத் தேவை:
\[
	\nicefrac{a}{b} = \nicefrac{c}{d} \text{ ஆக இருக்கும்போது, அப்போதும் அப்போது மட்டுமே, } a \times d = b \times c
\]
இதனைப் பெற, $\Int \times (\Int \setminus \{0_\Int\})$ மீது
ஒரு சமானத் தொடர்பைப் பின்வருமாறு வரையறுக்கிறோம்:
\[
	\tuple{ a, b }\Ratequiv \tuple{ c, d} \text{ ஆக இருக்கும்போது, அப்போதும் அப்போது மட்டுமே, }a \times d = b \times c
\]
இது ஒரு சமானத் தொடர்பு எனச் சரிபார்க்க வேண்டும்.
இது $\Intequiv$ இன் நிலையைப் பெரிதும் ஒத்துள்ளது;
இதனை ஒரு பயிற்சியாக விடுகிறோம்.

% SEGMENT OLP-0043-S02
\begin{prob}
$\Ratequiv$ என்பது ஒரு சமானத் தொடர்பு எனக் காட்டுக.
\end{prob}

% SEGMENT OLP-0043-S03
ஆனால் இது பின்வருமாறு கூற அனுமதிக்கிறது:
\begin{defn}
விகிதமுறு எண்கள் என்பவை, முழு ஜோடிகளுக்கு
(இரண்டாவது உறுப்பு பூச்சியமற்றதாக உள்ள ஜோடிகளுக்கு)
$\Ratequiv$ இன் கீழ் கிடைக்கும் சமான வகுப்புகள்.
அதாவது, $\Rat =
\equivclass{(\Int \times (\Int\setminus \{0_\Int\}))}{\Ratequiv}$.
\end{defn}

% SEGMENT OLP-0043-S04
முழுக்களைப் போலவே, சில அடிப்படைச் செயல்களையும் வரையறுக்க விரும்புகிறோம்.
$\equivrep{i,j}{\Ratequiv}$ என்பது $\tuple{i, j}$ ஐ ஓர்
!!a{element} ஆகக் கொண்ட, $\Ratequiv$ இன் கீழான சமான வகுப்பாக
இருக்கும்போது, பின்வருமாறு கூறுகிறோம்:
\begin{align*}
	\equivrep{a, b}{\Ratequiv} + \equivrep{c, d}{\Ratequiv} &= \equivrep{ad + bc,  bd}{\Ratequiv}\\
	\equivrep{a, b}{\Ratequiv} \times \equivrep{c, d}{\Ratequiv} &= \equivrep{a  c, b d}{\Ratequiv}.
\intertext{இந்த விகிதமுறு எண்கள் மீது $r \leq s$ என்பதை வரையறுக்க,
$r \le s$ ஆக இருக்கும்போது, அப்போதும் அப்போது மட்டுமே,
$s - r$ எதிர்மமற்றது என்ற உண்மையைப் பயன்படுத்துகிறோம்;
அதாவது, $s - r$ என்பதை $\nicefrac{i}{j}$ என்ற வடிவில் எழுதலாம்;
இங்கு $i$ எதிர்மமற்றதும் $j$ நேர்மமானதும் ஆகும்:}
	\equivrep{a, b}{\Ratequiv} \leq \equivrep{c, d}{\Ratequiv} &\text{ ஆக இருக்கும்போது, அப்போதும் அப்போது மட்டுமே, }
	\equivrep{c, d}{\Ratequiv} - \equivrep{a, b}{\Ratequiv} =
	\equivrep{i_\Int, j_\Int}{\Ratequiv}
\end{align*}
இங்கு ஏதோ ஒரு $i \in \Nat$, $0 \neq j \in \Nat$ உள்ளன.

இந்த வரையறைகள் \emph{செயல்பட வேண்டிய} விதத்தில் செயல்படுகின்றனவா
எனப் பிறகு சரிபார்க்க வேண்டும்; இதனை \olref[check]{sec} இல் தருகிறோம்.
அவை உண்மையில் அவ்வாறே செயல்படுகின்றன!{}
இறுதியாக, முழுக்களை விகிதமுறு எண்களாகக் \emph{கருதுவதற்கு}
ஒரு வழி தேவை; எனவே ஒவ்வொரு $i \in \Int$ க்கும்,
$i_\Rat = \equivrep{i, 1_\Int}{\Ratequiv}$ என நிர்ணயிக்கிறோம்.
இவை அனைத்தும் சரியாகச் செயல்படுகின்றன என்பதை மீண்டும்
\olref[check]{sec} இல் சரிபார்க்கிறோம்.

% SEGMENT OLP-0043-S05
\begin{prob}
எந்த $i, j \in \Int$ க்கும்,
$(i + j)_\Rat = i_\Rat+ j_\Rat$ மற்றும் $(i \times j)_\Rat =
i_\Rat \times j_\Rat$ மற்றும் $i \leq j \liff i_\Rat \leq j_\Rat$
எனக் காட்டுக.
\end{prob}

\end{document}
